\documentclass[11pt,a4paper]{article}
\usepackage[utf8]{inputenc}
\usepackage[T1]{fontenc}
\usepackage[ngerman]{babel}
\usepackage{helvet}
\renewcommand{\familydefault}{\sfdefault}
\usepackage[left=2cm, right=2cm, top=2cm, bottom=2cm]{geometry}
\usepackage{amsmath,amssymb}
\usepackage{tikz}
\usepackage{pgfplots}
\pgfplotsset{compat=1.18}
\usetikzlibrary{arrows.meta,positioning,decorations.pathreplacing,calc}
\usepackage{tcolorbox}
\tcbuselibrary{skins,breakable}
\usepackage{enumitem}
\usepackage{booktabs}
\usepackage{array}
\usepackage{colortbl}
\usepackage{fancyhdr}
\usepackage{hyperref}
\usepackage{multicol}

% Farben
\definecolor{physikblue}{RGB}{44,90,160}
\definecolor{lightblue}{RGB}{220,235,250}
\definecolor{merkegreen}{RGB}{34,139,94}
\definecolor{warnred}{RGB}{198,40,40}
\definecolor{tipporange}{RGB}{230,126,34}
\definecolor{lightgray}{RGB}{245,245,245}

% Boxen
\newtcolorbox{merksatz}[1][]{
  enhanced, breakable,
  colback=merkegreen!10, colframe=merkegreen,
  fonttitle=\bfseries, title={Merksatz},
  boxrule=1pt, arc=2pt, left=5pt, right=5pt, #1
}

\newtcolorbox{formelbox}[1][]{
  enhanced, breakable,
  colback=lightblue, colframe=physikblue,
  fonttitle=\bfseries, title={#1},
  boxrule=1pt, arc=2pt, left=5pt, right=5pt
}

\newtcolorbox{beispiel}[1][]{
  enhanced, breakable,
  colback=white, colframe=physikblue,
  fonttitle=\bfseries, title={Beispiel},
  boxrule=0.5pt, arc=2pt, left=5pt, right=5pt, #1
}

\newtcolorbox{warnung}{
  enhanced, breakable,
  colback=warnred!10, colframe=warnred,
  fonttitle=\bfseries, title={Häufiger Fehler},
  boxrule=1pt, arc=2pt, left=5pt, right=5pt
}

\newtcolorbox{tipp}{
  enhanced, breakable,
  colback=tipporange!10, colframe=tipporange,
  fonttitle=\bfseries, title={Tipp},
  boxrule=1pt, arc=2pt, left=5pt, right=5pt
}

\newtcolorbox{aufgabe}[1][]{
  enhanced, breakable,
  colback=lightgray, colframe=gray,
  fonttitle=\bfseries, title={#1},
  boxrule=0.5pt, arc=2pt, left=5pt, right=5pt
}

\definecolor{alltagpurple}{RGB}{106,27,154}
\newtcolorbox{alltag}{
  enhanced, breakable,
  colback=alltagpurple!8, colframe=alltagpurple,
  fonttitle=\bfseries, title={Was hat das mit mir zu tun?},
  boxrule=1pt, arc=2pt, left=5pt, right=5pt
}

% Header/Footer
\pagestyle{fancy}
\fancyhf{}
\fancyhead[L]{\small\color{gray}Kinematik -- Klasse 10E}
\fancyhead[R]{\small\color{gray}\thepage}
\renewcommand{\headrulewidth}{0.4pt}

% Section Styling
\usepackage{titlesec}
\titleformat{\section}{\Large\bfseries\color{physikblue}}{\thesection.}{0.5em}{}[\vspace{-0.3em}\color{physikblue}\titlerule[1pt]]
\titleformat{\subsection}{\large\bfseries\color{physikblue}}{\thesubsection}{0.5em}{}

\setlist{itemsep=2pt, topsep=3pt}

\begin{document}

% === KOPFZEILE ===
\begin{center}
{\Huge\bfseries\color{physikblue} Bewegungslehre -- Kinematik}\\[0.3cm]
{\large Klausurvorbereitung}\\[0.2cm]
{\large Klasse 10E}
\end{center}

\vspace{0.5cm}

\begin{tcolorbox}[colback=lightblue, colframe=physikblue, boxrule=1pt]
\textbf{Themen:} Gleichförmige Bewegung $\bullet$ Beschleunigte Bewegung $\bullet$ Zusammengesetzte Bewegungen $\bullet$ Freier Fall $\bullet$ Luftwiderstand \& Grenzgeschwindigkeit \textit{(nur qualitativ)}
\hfill \textbf{Merke:} $g \approx 10\,\mathrm{m/s^2}$
\end{tcolorbox}

\vspace{0.5cm}

\begin{tcolorbox}[colback=tipporange!10, colframe=tipporange, boxrule=1pt, title=\textbf{Vorwissens-Check: Bin ich bereit?}]
Bevor du startest -- kannst du das noch?

\begin{tabular}{|p{7cm}|c|c|p{4cm}|}
\hline
\rowcolor{tipporange!20}
\textbf{Vorwissen} & \textbf{Ja} & \textbf{Nein} & \textbf{Falls Nein: Wiederholen} \\
\hline
Gleichungen umstellen & $\square$ & $\square$ & Mathe Kl. 8 \\
\hline
Quadratwurzeln berechnen & $\square$ & $\square$ & Mathe Kl. 9 \\
\hline
SI-Einheiten (m, s, kg) & $\square$ & $\square$ & Physik Kl. 7 \\
\hline
Diagramme ablesen & $\square$ & $\square$ & Mathe Kl. 7 \\
\hline
\end{tabular}

\vspace{0.3cm}
\textbf{Schnelltest:} Rechne $\sqrt{25} = $ \underline{\hspace{1cm}} \quad und \quad $\dfrac{100}{4} = $ \underline{\hspace{1cm}}
\hfill {\small (Lösungen: 5 und 25)}
\end{tcolorbox}

\vspace{0.3cm}
\tableofcontents
\newpage

% ============================================================================
% KAPITEL 0: ÜBERBLICK BEWEGUNGEN
% ============================================================================
\section*{Vorab: Grundbegriffe der Bewegungslehre}
\addcontentsline{toc}{section}{Vorab: Grundbegriffe der Bewegungslehre}

\begin{merksatz}
\textbf{Bewegung} ist die \textbf{Änderung des Ortes} eines Körpers bezogen auf einen Bezugspunkt.

\textbf{Ruhe} bedeutet, dass sich der Ort eines Körpers \textbf{nicht ändert} -- er bleibt relativ zum Bezugspunkt an derselben Stelle.

\textit{Wichtig: Bewegung und Ruhe sind immer relativ! Ein Fahrgast im Zug ruht relativ zum Zug, bewegt sich aber relativ zur Landschaft.}
\end{merksatz}

\vspace{0.5cm}
Jede Bewegung lässt sich nach \textbf{zwei Kriterien} einteilen:

\vspace{0.5cm}

\begin{minipage}[t]{0.48\textwidth}
\begin{tcolorbox}[colback=physikblue!10, colframe=physikblue, title=\textbf{1. Bahnform}]
\begin{itemize}[leftmargin=0.5cm]
\item \textbf{geradlinig}\\
{\small Bewegung auf einer Geraden}\\
{\small (Auto auf Autobahn, freier Fall)}
\item \textbf{krummlinig}\\
{\small Bewegung auf gekrümmter Bahn}\\
{\small (Achterbahn, Slalom, Wurf)}\\[0.3em]
\textit{Spezialfall:} \textbf{kreisförmig}\\
{\small (Karussell, Erde um Sonne)}
\end{itemize}
\end{tcolorbox}
\end{minipage}
\hfill
\begin{minipage}[t]{0.48\textwidth}
\begin{tcolorbox}[colback=merkegreen!10, colframe=merkegreen, title=\textbf{2. Bewegungsart}]
\begin{itemize}[leftmargin=0.5cm]
\item \textbf{gleichförmig}\\
{\small Geschwindigkeit $v$ bleibt konstant}\\
{\small (Tempomat, Zug auf freier Strecke)}
\item \textbf{gleichmäßig beschleunigt}\\
{\small Beschleunigung $a$ bleibt konstant}\\
{\small (Anfahren, Bremsen, freier Fall)}
\end{itemize}
\end{tcolorbox}
\end{minipage}

\vspace{0.5cm}

\begin{center}
\textbf{Kombinationen:}\\[0.3cm]
\begin{tabular}{|l|c|c|}
\hline
\rowcolor{lightblue}
& \textbf{gleichförmig} & \textbf{gleichmäßig beschleunigt} \\
\hline
\textbf{geradlinig} & Auto mit Tempomat & Bremsen, freier Fall \\
\hline
\textbf{krummlinig} & Karussell & Waagerechter Wurf \\
\hline
\end{tabular}
\end{center}

\vspace{0.5cm}

\begin{tcolorbox}[colback=physikblue!5, colframe=physikblue, boxrule=1pt, title=\textbf{Kinematik vs. Dynamik -- Was ist der Unterschied?}]
\textbf{Kinematik} beschreibt \textbf{WIE} sich etwas bewegt:
\begin{itemize}
\item Geschwindigkeit, Beschleunigung, Weg, Zeit
\item Keine Kräfte, keine Ursachen
\item Fragen: \textit{Wie schnell? Wie weit? Wie lange?}
\end{itemize}

\textbf{Dynamik} (später) erklärt \textbf{WARUM} sich etwas bewegt:
\begin{itemize}
\item Kräfte als Ursache für Bewegungsänderungen
\item Newton'sche Gesetze
\item Fragen: \textit{Welche Kraft? Warum beschleunigt es?}
\end{itemize}

\textbf{In dieser Klausur:} Wir bleiben bei der \textbf{Kinematik} -- die Beschreibung der Bewegung.
\end{tcolorbox}

\newpage

% ============================================================================
% KAPITEL 1: GLEICHFÖRMIGE BEWEGUNG
% ============================================================================
\section{Gleichförmige Bewegung}

\begin{tcolorbox}[colback=merkegreen!5, colframe=merkegreen, boxrule=0.5pt]
\textbf{Nach diesem Kapitel kannst du...}
\begin{itemize}[itemsep=1pt]
\item[$\checkmark$] Geschwindigkeiten zwischen km/h und m/s umrechnen
\item[$\checkmark$] mit der Formel $v = s/t$ arbeiten und nach allen Größen umstellen
\item[$\checkmark$] $s(t)$- und $v(t)$-Diagramme zeichnen und interpretieren
\item[$\checkmark$] die Steigung im $s(t)$-Diagramm als Geschwindigkeit deuten
\end{itemize}
\end{tcolorbox}

\begin{alltag}
\textbf{Google Maps, Bahn-App, Navi} -- alle berechnen Ankunftszeiten mit genau dieser Formel!\\[0.3em]
Wenn du fragst „Wie lange brauche ich noch?" rechnet dein Handy: $t = \frac{s}{v}$.\\[0.3em]
\textbf{Führerschein:} Der Reaktionsweg wird mit $v \cdot t_{\text{Reaktion}}$ berechnet -- gleichförmige Bewegung!
\end{alltag}

\subsection{Definition und Grundformel}

Eine Bewegung heißt \textbf{gleichförmig}, wenn die Geschwindigkeit \textbf{konstant} bleibt -- der Körper legt in gleichen Zeitabschnitten gleiche Strecken zurück.

\begin{formelbox}[Grundformel der gleichförmigen Bewegung]
\[
\boxed{v = \frac{s}{t} = \frac{\Delta s}{\Delta t}}
\]
\begin{center}
\begin{tabular}{lcl}
$v$ & = & Geschwindigkeit in m/s \\
$s$ & = & zurückgelegte Strecke in m \\
$t$ & = & benötigte Zeit in s
\end{tabular}
\end{center}
\end{formelbox}

\begin{merksatz}
Die Geschwindigkeit gibt an, welche Strecke ein Körper \textbf{pro Zeiteinheit} zurücklegt.
\end{merksatz}

\subsection{Einheitenumrechnung}

\begin{tipp}
\textbf{km/h $\leftrightarrow$ m/s umrechnen:}
\begin{itemize}
\item km/h $\rightarrow$ m/s: \textbf{durch 3,6 teilen}
\item m/s $\rightarrow$ km/h: \textbf{mal 3,6 nehmen}
\end{itemize}

\vspace{0.3cm}
\textbf{Kopfrechenbare Werte (diese Werte parat haben!):}
\begin{center}
\begin{tabular}{|c|c|c|c|c|}
\hline
\rowcolor{lightblue}
\textbf{km/h} & 36 & 72 & 90 & 108 \\
\hline
\textbf{m/s} & 10 & 20 & 25 & 30 \\
\hline
\end{tabular}
\end{center}

\textit{Wusstest du? 50 km/h in der Stadt entspricht etwa 14 m/s -- pro Sekunde legst du also 14 Meter zurück. In der Reaktionszeit (ca. 1 s) fährst du blind 14 m weiter!}
\end{tipp}

\subsection{Weg-Zeit-Diagramm (s-t-Diagramm)}

Im $s(t)$-Diagramm wird die \textbf{Strecke} auf der y-Achse und die \textbf{Zeit} auf der x-Achse aufgetragen.

\begin{center}
\begin{tikzpicture}
\begin{axis}[
    width=10cm, height=6cm,
    xlabel={$t$ in s},
    ylabel={$s$ in m},
    xmin=0, xmax=6,
    ymin=0, ymax=60,
    grid=major,
    axis lines=left,
    legend pos=north west,
    legend style={font=\small}
]
\addplot[thick, physikblue, domain=0:5] {10*x};
\addplot[thick, merkegreen, domain=0:5] {5*x};
\addplot[thick, warnred, dashed, domain=0:5] {30};
\legend{$v = 10$ m/s (schnell), $v = 5$ m/s (langsam), Stillstand ($v = 0$)}
\end{axis}
\end{tikzpicture}
\end{center}

\begin{merksatz}
Im $s(t)$-Diagramm gilt:
\begin{itemize}
\item \textbf{Steigung = Geschwindigkeit} $\quad v = \dfrac{\Delta s}{\Delta t}$
\item Steile Gerade $\rightarrow$ hohe Geschwindigkeit
\item Flache Gerade $\rightarrow$ niedrige Geschwindigkeit
\item Horizontale Linie $\rightarrow$ Stillstand ($v = 0$)
\end{itemize}
\end{merksatz}

\subsection{Geschwindigkeit-Zeit-Diagramm (v-t-Diagramm)}

Im $v(t)$-Diagramm wird die \textbf{Geschwindigkeit} auf der y-Achse aufgetragen.

\begin{center}
\begin{tikzpicture}
\begin{axis}[
    width=10cm, height=5cm,
    xlabel={$t$ in s},
    ylabel={$v$ in m/s},
    xmin=0, xmax=6,
    ymin=0, ymax=15,
    grid=major,
    axis lines=left
]
\addplot[thick, physikblue, domain=0:5] {10};
\addlegendentry{$v = 10$ m/s (konstant)}
\fill[physikblue!30] (0,0) rectangle (5,10);
\node at (2.5,5) {\small Fläche = Weg};
\end{axis}
\end{tikzpicture}
\end{center}

\begin{merksatz}
Im $v(t)$-Diagramm gilt:
\begin{itemize}
\item \textbf{Horizontale Linie} = gleichförmige Bewegung (konstante Geschwindigkeit)
\item \textbf{Steigung = Beschleunigung} $\quad a = \dfrac{\Delta v}{\Delta t}$
\item \textbf{Fläche unter dem Graphen = zurückgelegter Weg}
\end{itemize}
\end{merksatz}

\begin{beispiel}
Ein Auto fährt mit konstanter Geschwindigkeit 72 km/h. Welchen Weg legt es in 25 Sekunden zurück?

\textbf{Gegeben:} $v = 72\,\mathrm{km/h}$, \quad $t = 25\,\mathrm{s}$

\textbf{Gesucht:} $s$

\textbf{Formel:} $s = v \cdot t$

\textbf{Einheiten umrechnen:} $v = \dfrac{72\,\mathrm{km/h}}{3{,}6} = 20\,\dfrac{\mathrm{m}}{\mathrm{s}}$

\textbf{Einsetzen:}
\[
s = v \cdot t = 20\,\frac{\mathrm{m}}{\mathrm{s}} \cdot 25\,\mathrm{s}
\]

\textbf{Rechnung:}
\[
s = 20 \cdot 25\,\frac{\mathrm{m} \cdot \mathrm{s}}{\mathrm{s}} = 500\,\mathrm{m}
\]

\textbf{Ergebnis:} $\boxed{s = 500\,\mathrm{m}}$
\end{beispiel}

% ============================================================================
% KAPITEL 2: BESCHLEUNIGTE BEWEGUNG
% ============================================================================
\newpage
\section{Gleichmäßig beschleunigte Bewegung}

\begin{tcolorbox}[colback=merkegreen!5, colframe=merkegreen, boxrule=0.5pt]
\textbf{Nach diesem Kapitel kannst du...}
\begin{itemize}[itemsep=1pt]
\item[$\checkmark$] die Beschleunigung als Geschwindigkeitsänderung pro Zeit verstehen
\item[$\checkmark$] die Formeln $v(t) = at + v_0$ und $s(t) = \frac{1}{2}at^2 + v_0 t$ anwenden
\item[$\checkmark$] den Zusammenhang zwischen $v(t)$- und $s(t)$-Diagramm erkennen
\item[$\checkmark$] die Fläche unter dem $v(t)$-Graphen als Weg interpretieren
\end{itemize}
\end{tcolorbox}

\begin{alltag}
\textbf{Führerschein-Pflicht:} Der \textbf{Bremsweg} ist eine der wichtigsten Prüfungsfragen! \\[0.3em]
Bei doppelter Geschwindigkeit: \textbf{4-facher Bremsweg} -- das rettet Leben.\\[0.3em]
\textbf{0-100 in X Sekunden} -- Autotests geben immer die Beschleunigung an. Schnelle E-Autos schaffen $a \approx 4\,\mathrm{m/s^2}$!
\end{alltag}

\subsection{Definition der Beschleunigung}

Wenn sich die Geschwindigkeit ändert, liegt eine \textbf{beschleunigte Bewegung} vor.

\begin{formelbox}[Beschleunigung]
\[
\boxed{a = \frac{\Delta v}{\Delta t} = \frac{v - v_0}{t}}
\]
\begin{center}
\begin{tabular}{lcl}
$a$ & = & Beschleunigung in m/s² \\
$\Delta v$ & = & Geschwindigkeitsänderung in m/s \\
$\Delta t$ & = & Zeitintervall in s
\end{tabular}
\end{center}
\end{formelbox}

\begin{merksatz}
\begin{itemize}
\item $a > 0$: Geschwindigkeit nimmt zu (Beschleunigen)
\item $a < 0$: Geschwindigkeit nimmt ab (Bremsen/Verzögern)
\item $a = 0$: Geschwindigkeit bleibt konstant (gleichförmig)
\end{itemize}
\end{merksatz}

\subsection{Die drei Grundformeln}

\begin{formelbox}[Formeln der gleichmäßig beschleunigten Bewegung]
\begin{align}
v(t) &= a \cdot t + v_0 \tag{Geschwindigkeit}\\[0.3em]
s(t) &= \frac{1}{2} a \cdot t^2 + v_0 \cdot t + s_0 \tag{Weg}\\[0.3em]
a(t) &= \text{konstant} \tag{Beschleunigung}
\end{align}

\vspace{0.3cm}
\textbf{Ohne Anfangsbedingungen} (Start aus Ruhe bei $s_0 = 0$):
\begin{align}
v(t) &= a \cdot t \\
s(t) &= \frac{1}{2} a \cdot t^2
\end{align}

\textbf{Nützliche Zusatzformel} (ohne Zeit -- auch \textbf{„zeitlose Formel"}):
\[
v^2 = v_0^2 + 2 \cdot a \cdot s
\]
\textit{Für Interessierte: Die Herleitung findest du im Anhang \ref{sec:herleitung}.}
\end{formelbox}

\subsection{Diagramme der beschleunigten Bewegung}

\begin{center}
\begin{tikzpicture}
% v-t-Diagramm
\begin{axis}[
    name=vt,
    width=7cm, height=5cm,
    xlabel={$t$ in s},
    ylabel={$v$ in m/s},
    xmin=0, xmax=5,
    ymin=0, ymax=25,
    grid=major,
    axis lines=left,
    title={\textbf{v-t-Diagramm}}
]
\addplot[thick, physikblue, domain=0:5] {5*x};
\fill[physikblue!30] (0,0) -- (4,0) -- (4,20) -- cycle;
\node at (2,5) {\small Fläche = $s$};
\end{axis}

% s-t-Diagramm
\begin{axis}[
    at={(8cm,0)},
    width=7cm, height=5cm,
    xlabel={$t$ in s},
    ylabel={$s$ in m},
    xmin=0, xmax=5,
    ymin=0, ymax=70,
    grid=major,
    axis lines=left,
    title={\textbf{s-t-Diagramm}}
]
\addplot[thick, merkegreen, domain=0:5] {2.5*x^2};
\end{axis}
\end{tikzpicture}
\end{center}

\begin{merksatz}
Bei gleichmäßig beschleunigter Bewegung (aus Ruhe):
\begin{itemize}
\item $v(t)$-Diagramm: \textbf{Gerade durch den Ursprung} (Steigung = $a$)
\item $s(t)$-Diagramm: \textbf{Parabel} (nach oben geöffnet)
\item Fläche unter $v(t)$ = zurückgelegter Weg
\end{itemize}
\end{merksatz}

\begin{formelbox}[Diagramm-Interpretation: Steigung und Fläche]
\begin{center}
\renewcommand{\arraystretch}{1.5}
\begin{tabular}{|l|c|c|}
\hline
\rowcolor{lightblue}
\textbf{Diagramm} & \textbf{Steigung} & \textbf{Fläche} \\
\hline
$s(t)$-Diagramm & Geschwindigkeit $v$ & --- \\
\hline
$v(t)$-Diagramm & Beschleunigung $a$ & Weg $\Delta s$ \\
\hline
$a(t)$-Diagramm & --- & Geschwindigkeitsänderung $\Delta v$ \\
\hline
\end{tabular}
\end{center}

\vspace{0.3cm}
\textbf{Merkregel:} Die Fläche unter einem Graphen ergibt die darüberliegende Größe:
\begin{itemize}
\item Fläche unter $a(t)$ $\rightarrow$ $\Delta v$ (Geschwindigkeitsänderung)
\item Fläche unter $v(t)$ $\rightarrow$ $\Delta s$ (Weg)
\end{itemize}
\end{formelbox}

\begin{beispiel}
Ein Auto beschleunigt aus dem Stand ($v_0 = 0$) mit $a = 4\,\mathrm{m/s^2}$.

\textbf{a)} Welche Geschwindigkeit hat es nach 5 Sekunden?

\textbf{Gegeben:} $a = 4\,\mathrm{m/s^2}$, \quad $t = 5\,\mathrm{s}$, \quad $v_0 = 0$

\textbf{Formel:} $v = a \cdot t + v_0 = a \cdot t$ \quad (da $v_0 = 0$)

\textbf{Einsetzen:}
\[
v = 4\,\frac{\mathrm{m}}{\mathrm{s}^2} \cdot 5\,\mathrm{s}
\]

\textbf{Rechnung:}
\[
v = 4 \cdot 5\,\frac{\mathrm{m} \cdot \mathrm{s}}{\mathrm{s}^2} = 20\,\frac{\mathrm{m}}{\mathrm{s}}
\]

\textbf{Ergebnis:} $\boxed{v = 20\,\mathrm{m/s} = 72\,\mathrm{km/h}}$

\vspace{0.3cm}
\textbf{b)} Welchen Weg legt es in dieser Zeit zurück?

\textbf{Formel:} $s = \dfrac{1}{2} a \cdot t^2$ \quad (da $v_0 = 0$, $s_0 = 0$)

\textbf{Einsetzen:}
\[
s = \frac{1}{2} \cdot 4\,\frac{\mathrm{m}}{\mathrm{s}^2} \cdot (5\,\mathrm{s})^2
\]

\textbf{Rechnung:}
\[
s = \frac{1}{2} \cdot 4 \cdot 25\,\frac{\mathrm{m} \cdot \mathrm{s}^2}{\mathrm{s}^2} = 2 \cdot 25\,\mathrm{m} = 50\,\mathrm{m}
\]

\textbf{Ergebnis:} $\boxed{s = 50\,\mathrm{m}}$
\end{beispiel}

% ============================================================================
% KAPITEL 3: ZUSAMMENGESETZTE BEWEGUNGEN
% ============================================================================
\newpage
\section{Zusammengesetzte Bewegungen}

\begin{tcolorbox}[colback=merkegreen!5, colframe=merkegreen, boxrule=0.5pt]
\textbf{Nach diesem Kapitel kannst du...}
\begin{itemize}[itemsep=1pt]
\item[$\checkmark$] Bremsvorgänge berechnen (Bremszeit, Bremsweg)
\item[$\checkmark$] die zeitlose Formel $v^2 = v_0^2 + 2as$ anwenden
\item[$\checkmark$] Überholaufgaben durch Gleichsetzen von Bewegungsgleichungen lösen
\item[$\checkmark$] Vorzeichen bei entgegengesetzten Bewegungen richtig setzen
\end{itemize}
\end{tcolorbox}

\begin{alltag}
\textbf{Überholen auf der Landstraße} -- lebensgefährlich, wenn man falsch rechnet!\\[0.3em]
Wie viel Strecke brauche ich wirklich? Reicht die Lücke vor dem Gegenverkehr?\\[0.3em]
\textbf{Zugfahrplan:} Wann treffen sich zwei Züge? Wann überholt der ICE den Regionalzug?\\[0.3em]
Diese Berechnungen retten Leben und verhindern Unfälle.
\end{alltag}

\subsection{Prinzip}

Viele reale Bewegungen bestehen aus \textbf{mehreren Phasen} mit unterschiedlichen Bewegungsarten.

\begin{merksatz}
Bei zusammengesetzten Bewegungen:
\begin{enumerate}
\item \textbf{Identifiziere} die einzelnen Bewegungsphasen
\item \textbf{Berechne} jede Phase separat
\item \textbf{Beachte:} Endwerte einer Phase = Anfangswerte der nächsten
\end{enumerate}
\end{merksatz}

\subsection{Bremsvorgang}

Beim Bremsen gilt: $a < 0$ (negative Beschleunigung = Verzögerung)

\begin{formelbox}[Bremsformeln]
\textbf{Bremszeit} (Zeit bis zum Stillstand):
\[
t_{\text{brems}} = \frac{v_0}{|a|}
\]

\textbf{Bremsweg:}
\[
s_{\text{brems}} = \frac{v_0^2}{2 \cdot |a|}
\]
\end{formelbox}

\begin{warnung}
\textbf{Lebenswichtig:} Der Bremsweg hängt von $v^2$ ab!
\begin{center}
\begin{tabular}{|c|c|c|}
\hline
\rowcolor{lightblue}
\textbf{Tempo} & \textbf{Bremsweg} & \textbf{Vergleich} \\
\hline
30 km/h & ca. 4 m & --- \\
\hline
50 km/h & ca. 12 m & 3× so weit \\
\hline
100 km/h & ca. 50 m & 12× so weit! \\
\hline
\end{tabular}
\end{center}
\textit{Deshalb gilt in Spielstraßen Tempo 30 -- damit Kinder noch eine Chance haben.}
\end{warnung}

\begin{beispiel}
Ein Auto fährt mit $v_0 = 36\,\mathrm{km/h}$. Bei Wegmarke $s_0 = 500\,\mathrm{m}$ beginnt der Fahrer zu bremsen mit $a = -2\,\mathrm{m/s^2}$.

\textbf{Gegeben:} $v_0 = 36\,\mathrm{km/h}$, \quad $s_0 = 500\,\mathrm{m}$, \quad $a = -2\,\mathrm{m/s^2}$

\textbf{Einheiten umrechnen:} $v_0 = \dfrac{36}{3{,}6}\,\dfrac{\mathrm{m}}{\mathrm{s}} = 10\,\dfrac{\mathrm{m}}{\mathrm{s}}$

\vspace{0.3cm}
\textbf{a) Bremszeit:}

\textbf{Formel:} $t = \dfrac{v_0}{|a|}$

\textbf{Einsetzen:}
\[
t = \frac{10\,\frac{\mathrm{m}}{\mathrm{s}}}{2\,\frac{\mathrm{m}}{\mathrm{s}^2}}
\]

\textbf{Rechnung:}
\[
t = \frac{10}{2}\,\frac{\mathrm{m} \cdot \mathrm{s}^2}{\mathrm{s} \cdot \mathrm{m}} = 5\,\mathrm{s}
\]

\textbf{Ergebnis:} $\boxed{t = 5\,\mathrm{s}}$

\vspace{0.3cm}
\textbf{b) Bremsweg:}

\textbf{Formel:} $s_{\text{brems}} = \dfrac{v_0^2}{2 \cdot |a|}$

\textbf{Einsetzen:}
\[
s_{\text{brems}} = \frac{\left(10\,\frac{\mathrm{m}}{\mathrm{s}}\right)^2}{2 \cdot 2\,\frac{\mathrm{m}}{\mathrm{s}^2}}
\]

\textbf{Rechnung:}
\[
s_{\text{brems}} = \frac{100\,\frac{\mathrm{m}^2}{\mathrm{s}^2}}{4\,\frac{\mathrm{m}}{\mathrm{s}^2}} = \frac{100}{4}\,\frac{\mathrm{m}^2 \cdot \mathrm{s}^2}{\mathrm{s}^2 \cdot \mathrm{m}} = 25\,\mathrm{m}
\]

\textbf{Ergebnis:} $\boxed{s_{\text{brems}} = 25\,\mathrm{m}}$

\vspace{0.5cm}
\begin{tcolorbox}[colback=physikblue!5, colframe=physikblue, boxrule=0.5pt, title=\textbf{Variante B: Mit zeitloser Formel}]
Alternativ kann der Bremsweg \textbf{ohne Umweg über die Zeit} berechnet werden:

\textbf{Zeitlose Formel:} $v^2 = v_0^2 + 2 \cdot a \cdot s$

\textbf{Umstellen} (mit $v = 0$ für Stillstand):
\[
0 = v_0^2 + 2 \cdot a \cdot s \quad \Rightarrow \quad s = -\frac{v_0^2}{2 \cdot a} = \frac{v_0^2}{2 \cdot |a|}
\]

\textbf{Einsetzen:}
\[
s = \frac{(10\,\mathrm{m/s})^2}{2 \cdot 2\,\mathrm{m/s}^2} = \frac{100}{4}\,\mathrm{m} = 25\,\mathrm{m} \quad \checkmark
\]

\textit{Vorteil: Kein Zwischenschritt über $t$ nötig!}\\[0.3em]
\textit{(Warum funktioniert das? $\rightarrow$ Anhang \ref{sec:herleitung})}
\end{tcolorbox}

\vspace{0.3cm}
\textbf{c) Endposition:}

\textbf{Formel:} $s_{\text{Ende}} = s_0 + s_{\text{brems}}$

\textbf{Einsetzen und Rechnung:}
\[
s_{\text{Ende}} = 500\,\mathrm{m} + 25\,\mathrm{m} = 525\,\mathrm{m}
\]

\textbf{Ergebnis:} $\boxed{s_{\text{Ende}} = 525\,\mathrm{m}}$
\end{beispiel}

\subsection{Bewegung mit Anfangsbedingungen}

\begin{formelbox}[Allgemeine Weg-Zeit-Gleichung]
\[
\boxed{s(t) = \frac{1}{2} a \cdot t^2 + v_0 \cdot t + s_0}
\]

Diese Formel vereint alle Fälle:
\begin{itemize}
\item $a = 0$: gleichförmige Bewegung $\rightarrow s(t) = v_0 \cdot t + s_0$
\item $v_0 = 0$, $s_0 = 0$: Start aus Ruhe $\rightarrow s(t) = \frac{1}{2}at^2$
\end{itemize}
\end{formelbox}

\subsection{Überholmanöver}

Bei Überholaufgaben müssen zwei Bewegungsgleichungen \textbf{gleichgesetzt} werden.

\begin{beispiel}
Radfahrer A startet bei $s_0 = 10\,\mathrm{m}$ mit $v_A = 3\,\mathrm{m/s}$.\\
Radfahrer B startet bei $s_0 = 0$ mit $v_B = 5\,\mathrm{m/s}$.

\textbf{Wann und wo überholt B den A?}

\textbf{1. Bewegungsgleichungen aufstellen:}
\begin{align*}
s_A(t) &= v_A \cdot t + s_{0,A} = 3\,\frac{\mathrm{m}}{\mathrm{s}} \cdot t + 10\,\mathrm{m} \\
s_B(t) &= v_B \cdot t + s_{0,B} = 5\,\frac{\mathrm{m}}{\mathrm{s}} \cdot t + 0
\end{align*}

\textbf{2. Gleichsetzen (Bedingung: B überholt A, wenn $s_A = s_B$):}
\[
s_A(t) = s_B(t)
\]

\textbf{3. Umstellen:}
\begin{align*}
v_A \cdot t + s_{0,A} &= v_B \cdot t \\
s_{0,A} &= v_B \cdot t - v_A \cdot t \\
s_{0,A} &= t \cdot (v_B - v_A) \\
t &= \frac{s_{0,A}}{v_B - v_A}
\end{align*}

\textbf{4. Einsetzen:}
\[
t = \frac{10\,\mathrm{m}}{5\,\frac{\mathrm{m}}{\mathrm{s}} - 3\,\frac{\mathrm{m}}{\mathrm{s}}} = \frac{10\,\mathrm{m}}{2\,\frac{\mathrm{m}}{\mathrm{s}}} = \boxed{5\,\mathrm{s}}
\]

\textbf{5. Position berechnen:}
\[
s = v_B \cdot t = 5\,\frac{\mathrm{m}}{\mathrm{s}} \cdot 5\,\mathrm{s} = \boxed{25\,\mathrm{m}}
\]
\end{beispiel}

\begin{beispiel}
\textbf{Entgegengesetzte Fahrt: Rostock $\leftrightarrow$ Graal-Müritz}

Ein Zug startet in Rostock ($s_0 = 0$) Richtung Graal-Müritz mit $v_1 = 20\,\mathrm{m/s}$.\\
Gleichzeitig startet ein Zug in Graal-Müritz ($s_0 = 25\,\mathrm{km}$) Richtung Rostock mit $v_2 = 30\,\mathrm{m/s}$.

\textbf{Wann und wo treffen sich die Züge?}

\begin{warnung}
\textbf{Vorsicht: Vorzeichen!}

Der zweite Zug fährt in \textbf{negativer Richtung} (nach Rostock = kleinere $s$-Werte).\\
Daher: $v_2 = -30\,\mathrm{m/s}$ (negativ!)
\end{warnung}

\textbf{Gegeben:}
\begin{itemize}[itemsep=1pt]
\item Zug 1: $s_{0,1} = 0$, \quad $v_1 = +20\,\mathrm{m/s}$
\item Zug 2: $s_{0,2} = 25\,\mathrm{km} = 25000\,\mathrm{m}$, \quad $v_2 = -30\,\mathrm{m/s}$
\end{itemize}

\textbf{1. Bewegungsgleichungen aufstellen:}
\begin{align*}
s_1(t) &= v_1 \cdot t + s_{0,1} = (+20)\,\frac{\mathrm{m}}{\mathrm{s}} \cdot t + 0 \\[0.3em]
s_2(t) &= v_2 \cdot t + s_{0,2} = (-30)\,\frac{\mathrm{m}}{\mathrm{s}} \cdot t + 25000\,\mathrm{m}
\end{align*}

\textbf{2. Gleichsetzen (Treffpunkt: $s_1 = s_2$):}
\[
20\,\frac{\mathrm{m}}{\mathrm{s}} \cdot t = -30\,\frac{\mathrm{m}}{\mathrm{s}} \cdot t + 25000\,\mathrm{m}
\]

\textbf{3. Umstellen:}
\begin{align*}
20\,\frac{\mathrm{m}}{\mathrm{s}} \cdot t + 30\,\frac{\mathrm{m}}{\mathrm{s}} \cdot t &= 25000\,\mathrm{m} \\[0.3em]
50\,\frac{\mathrm{m}}{\mathrm{s}} \cdot t &= 25000\,\mathrm{m} \\[0.3em]
t &= \frac{25000\,\mathrm{m}}{50\,\frac{\mathrm{m}}{\mathrm{s}}}
\end{align*}

\textbf{4. Rechnung:}
\[
t = \frac{25000}{50}\,\frac{\mathrm{m} \cdot \mathrm{s}}{\mathrm{m}} = 500\,\mathrm{s} = 8\,\mathrm{min}\,20\,\mathrm{s}
\]

\textbf{Ergebnis Treffzeit:} $\boxed{t = 500\,\mathrm{s}}$

\textbf{5. Treffpunkt berechnen:}
\[
s = 20\,\frac{\mathrm{m}}{\mathrm{s}} \cdot 500\,\mathrm{s} = 10000\,\mathrm{m} = 10\,\mathrm{km}
\]

\textbf{Ergebnis Treffpunkt:} $\boxed{s = 10\,\mathrm{km}}$ (von Rostock aus gemessen)

\begin{tipp}
\textbf{Merkhilfe für entgegengesetzte Bewegung:}

Die \textbf{Relativgeschwindigkeit} ist die \textbf{Summe} der Beträge: $|v_1| + |v_2| = 20 + 30 = 50\,\mathrm{m/s}$.

Schnellcheck: $t = \dfrac{25000\,\mathrm{m}}{50\,\mathrm{m/s}} = 500\,\mathrm{s}$ \checkmark
\end{tipp}
\end{beispiel}

\begin{warnung}
\textbf{Anfangsbedingungen und Vorzeichen nicht vergessen!}

Wenn ein Körper nicht bei $s = 0$ startet oder bereits eine Anfangsgeschwindigkeit hat, müssen $s_0$ und $v_0$ in die Formel eingesetzt werden!
\end{warnung}

% ============================================================================
% KAPITEL 4: FREIER FALL
% ============================================================================
\newpage
\section{Der freie Fall}

\begin{tcolorbox}[colback=merkegreen!5, colframe=merkegreen, boxrule=0.5pt]
\textbf{Nach diesem Kapitel kannst du...}
\begin{itemize}[itemsep=1pt]
\item[$\checkmark$] den freien Fall als gleichmäßig beschleunigte Bewegung mit $g \approx 10\,\mathrm{m/s^2}$ erkennen
\item[$\checkmark$] Fallzeit, Fallstrecke und Aufprallgeschwindigkeit berechnen
\item[$\checkmark$] den senkrechten Wurf nach oben analysieren
\item[$\checkmark$] die vereinfachte Formel $s = 5t^2$ kopfrechnerisch nutzen
\end{itemize}
\end{tcolorbox}

\begin{alltag}
\textbf{Handy fallen lassen?} Nach 0,5 Sekunden knallt es auf den Boden -- aus nur 1,25 m Höhe!\\[0.3em]
\textbf{Bungee-Jumping, Fallschirmspringen, Achterbahn} -- der freie Fall ist pure Physik.\\[0.3em]
\textbf{Galileo Galilei} widerlegte vor 400 Jahren den Irrtum, schwere Dinge fielen schneller.\\[0.3em]
Fun Fact: Auf dem Mond ließ ein Astronaut Hammer und Feder gleichzeitig fallen -- sie landeten zusammen!
\end{alltag}

\subsection{Definition}

\begin{merksatz}
Der \textbf{freie Fall} ist eine gleichmäßig beschleunigte Bewegung nach unten, bei der nur die \textbf{Gewichtskraft} wirkt (kein Luftwiderstand).

Die Beschleunigung heißt \textbf{Fallbeschleunigung}:
\[
g = 9{,}81\,\frac{\mathrm{m}}{\mathrm{s}^2} \approx \boxed{10\,\frac{\mathrm{m}}{\mathrm{s}^2}}
\]
\end{merksatz}

\begin{tipp}
In der Klausur \textbf{immer} $g \approx 10\,\mathrm{m/s^2}$ verwenden -- das macht Kopfrechnen möglich!
\end{tipp}

\subsection{Formeln des freien Falls}

\begin{formelbox}[Formeln des freien Falls]
\begin{center}
\renewcommand{\arraystretch}{1.8}
\begin{tabular}{|l|c|c|}
\hline
\rowcolor{lightblue}
\textbf{Größe} & \textbf{Exakt} & \textbf{Vereinfacht} \\
\hline
Fallgeschwindigkeit & $v = g \cdot t$ & $v = 10 \cdot t$ \\
\hline
Fallstrecke & $s = \frac{1}{2} g \cdot t^2$ & $s = 5 \cdot t^2$ \\
\hline
Fallzeit (aus $s$) & $t = \sqrt{\frac{2s}{g}}$ & $t = \sqrt{\frac{s}{5}}$ \\
\hline
\end{tabular}
\end{center}
\end{formelbox}

\subsection{Wertetabelle}

\begin{center}
\begin{tabular}{|c|c|c|c|c|c|c|}
\hline
\rowcolor{lightblue}
$t$ in s & 0 & 1 & 2 & 3 & 4 & 5 \\
\hline
$v$ in m/s & 0 & 10 & 20 & 30 & 40 & 50 \\
\hline
$s$ in m & 0 & 5 & 20 & 45 & 80 & 125 \\
\hline
\end{tabular}
\end{center}

\subsection{Diagramme des freien Falls}

\begin{center}
\begin{tikzpicture}
% v-t
\begin{axis}[
    name=vt,
    width=6.5cm, height=5cm,
    xlabel={$t$ in s},
    ylabel={$v$ in m/s},
    xmin=0, xmax=5,
    ymin=0, ymax=55,
    grid=major,
    axis lines=left,
    title={\small $v(t)$-Diagramm}
]
\addplot[thick, physikblue, domain=0:5] {10*x};
\end{axis}

% s-t
\begin{axis}[
    at={(7.5cm,0)},
    width=6.5cm, height=5cm,
    xlabel={$t$ in s},
    ylabel={$s$ in m},
    xmin=0, xmax=5,
    ymin=0, ymax=130,
    grid=major,
    axis lines=left,
    title={\small $s(t)$-Diagramm}
]
\addplot[thick, merkegreen, domain=0:5] {5*x^2};
\end{axis}
\end{tikzpicture}
\end{center}

\begin{beispiel}
Du stehst auf einer Brücke und lässt einen Stein fallen. Nach 3 Sekunden hörst du den Aufprall.

\textbf{Gegeben:} $t = 3\,\mathrm{s}$, \quad $g = 10\,\mathrm{m/s^2}$

\vspace{0.3cm}
\textbf{a) Höhe der Brücke:}

\textbf{Gesucht:} $s$ (Fallstrecke = Brückenhöhe)

\textbf{Formel:} $s = \dfrac{1}{2} g \cdot t^2 = 5\,\dfrac{\mathrm{m}}{\mathrm{s}^2} \cdot t^2$

\textbf{Einsetzen:}
\[
s = 5\,\frac{\mathrm{m}}{\mathrm{s}^2} \cdot (3\,\mathrm{s})^2
\]

\textbf{Rechnung:}
\[
s = 5\,\frac{\mathrm{m}}{\mathrm{s}^2} \cdot 9\,\mathrm{s}^2 = 5 \cdot 9\,\frac{\mathrm{m} \cdot \mathrm{s}^2}{\mathrm{s}^2} = 45\,\mathrm{m}
\]

\textbf{Ergebnis:} $\boxed{s = 45\,\mathrm{m}}$

\vspace{0.3cm}
\textbf{b) Aufprallgeschwindigkeit:}

\textbf{Gesucht:} $v$ (Geschwindigkeit beim Aufprall)

\textbf{Formel:} $v = g \cdot t$

\textbf{Einsetzen:}
\[
v = 10\,\frac{\mathrm{m}}{\mathrm{s}^2} \cdot 3\,\mathrm{s}
\]

\textbf{Rechnung:}
\[
v = 10 \cdot 3\,\frac{\mathrm{m} \cdot \mathrm{s}}{\mathrm{s}^2} = 30\,\frac{\mathrm{m}}{\mathrm{s}}
\]

\textbf{Ergebnis:} $\boxed{v = 30\,\mathrm{m/s} = 108\,\mathrm{km/h}}$
\end{beispiel}

\subsection{Senkrechter Wurf nach oben}

Beim Wurf nach oben gilt: Am höchsten Punkt ist $v = 0$.

\begin{formelbox}[Senkrechter Wurf nach oben]
\textbf{Steigzeit} (bis $v = 0$):
\[
t_{\text{steig}} = \frac{v_0}{g}
\]

\textbf{Steighöhe:}
\[
h_{\text{max}} = \frac{v_0^2}{2g}
\]

\textbf{Gesamtzeit} (hin und zurück):
\[
t_{\text{ges}} = 2 \cdot t_{\text{steig}} = \frac{2v_0}{g}
\]
\end{formelbox}

\begin{merksatz}
Die Rückfallgeschwindigkeit ist \textbf{gleich groß} wie die Abwurfgeschwindigkeit (nur entgegengesetzte Richtung).
\end{merksatz}

% ============================================================================
% KAPITEL 5: LUFTWIDERSTAND
% ============================================================================
\newpage
\section{Luftwiderstand und Grenzgeschwindigkeit}

\begin{tcolorbox}[colback=merkegreen!5, colframe=merkegreen, boxrule=0.5pt]
\textbf{Nach diesem Kapitel kannst du...}
\begin{itemize}[itemsep=1pt]
\item[$\checkmark$] erklären, warum $F_W$ bei höherer Geschwindigkeit stark zunimmt ($\sim v^2$)
\item[$\checkmark$] beschreiben, wie Grenzgeschwindigkeit entsteht ($F_G = F_W$)
\item[$\checkmark$] den Einfluss von Masse, Fläche und Luftdichte auf $v_{\text{grenz}}$ qualitativ begründen
\item[$\checkmark$] Alltagsphänomene (Fallschirm, Regentropfen) physikalisch erklären
\end{itemize}
\end{tcolorbox}

\begin{alltag}
\textbf{Warum tragen Radfahrer enge Anzüge?} Weniger Luftwiderstand = schneller!\\[0.3em]
\textbf{Warum öffnet sich der Fallschirm?} Große Fläche = viel Luftwiderstand = langsam und sicher landen.\\[0.3em]
\textbf{Felix Baumgartner} sprang aus 39 km Höhe und wurde schneller als der Schall -- weil dort oben kaum Luft ist!\\[0.3em]
\textbf{Regentropfen:} Ohne Luftwiderstand würden sie mit 200 km/h aufschlagen -- autsch!
\end{alltag}

\begin{tcolorbox}[colback=tipporange!20, colframe=tipporange, boxrule=2pt]
\centering
\textbf{\large Nur qualitatives Verständnis erforderlich!}\\[0.3em]
\normalsize
Für die Klausur: \textbf{Keine Berechnungen} zu Luftwiderstand oder Grenzgeschwindigkeit.\\
Ihr müsst nur die \textbf{Zusammenhänge verstehen und erklären} können.
\end{tcolorbox}

\vspace{0.5cm}

\subsection{Die Luftwiderstandskraft}

Wenn ein Körper durch Luft fällt, wirkt ihm die \textbf{Luftwiderstandskraft} entgegen.

\begin{formelbox}[Luftwiderstandskraft (nur qualitativ!)]
\[
F_W = \frac{1}{2} \cdot \rho \cdot c_W \cdot A \cdot v^2
\]

\textbf{Wichtig sind nur die Zusammenhänge:}
\begin{center}
\begin{tabular}{lcl}
$\rho$ (Luftdichte) größer & $\Rightarrow$ & $F_W$ größer \\
$A$ (Querschnitt) größer & $\Rightarrow$ & $F_W$ größer \\
$v$ (Geschwindigkeit) größer & $\Rightarrow$ & $F_W$ \textbf{viel} größer (quadratisch!)
\end{tabular}
\end{center}
\end{formelbox}

\begin{merksatz}
\textbf{Wichtig:} $F_W$ hängt von $v^2$ ab!
\begin{itemize}
\item Verdoppelt sich $v$ $\rightarrow$ $F_W$ vervierfacht sich!
\item Je schneller man fällt, desto größer wird der Luftwiderstand.
\end{itemize}
\end{merksatz}

\subsection{Der Widerstandsbeiwert $c_W$}

\begin{center}
\begin{tikzpicture}[scale=0.9]
% Platte
\draw[fill=gray!30] (0,0) rectangle (0.4,2);
\node[below,font=\small] at (0.2,-0.2) {Platte};
\node[below,font=\footnotesize] at (0.2,-0.6) {$c_W \approx 1{,}1$};

% Kugel
\draw[fill=gray!30] (4,1) circle (0.8);
\node[below,font=\small] at (4,-0.2) {Kugel};
\node[below,font=\footnotesize] at (4,-0.6) {$c_W \approx 0{,}45$};

% Tropfen
\draw[fill=gray!30] (8.2,1) ellipse (0.4 and 0.8);
\draw[fill=gray!30] (8.2,0.2) -- (9.4,1) -- (8.2,1.8) -- cycle;
\node[below,font=\small] at (8.6,-0.2) {Tropfenform};
\node[below,font=\footnotesize] at (8.6,-0.6) {$c_W \approx 0{,}04$};
\end{tikzpicture}
\end{center}

\begin{merksatz}
Je \textbf{stromlinienförmiger} ein Körper, desto kleiner sein $c_W$-Wert und desto geringer der Luftwiderstand.
\end{merksatz}

\subsection{Die Grenzgeschwindigkeit}

Beim Fall wirken zwei Kräfte:
\begin{itemize}
\item $F_G = m \cdot g$ (Gewichtskraft, nach unten)
\item $F_W = \frac{1}{2}\rho c_W A v^2$ (Luftwiderstand, nach oben)
\end{itemize}

\begin{center}
\begin{tikzpicture}
% Phase 1
\node at (0,2.5) {\textbf{Start}};
\fill[gray!50] (0,1) circle (0.4);
\draw[-{Stealth[length=3mm]},thick,warnred] (0,0.6) -- (0,-0.8) node[right] {$F_G$};
\draw[-{Stealth[length=2mm]},thick,physikblue] (0,1.4) -- (0,1.8) node[right] {\small $F_W$};
\node at (0,-1.5) {\small $F_G > F_W$};
\node at (0,-2) {\small $\Rightarrow$ beschleunigt};

% Phase 2
\node at (4,2.5) {\textbf{Später}};
\fill[gray!50] (4,1) circle (0.4);
\draw[-{Stealth[length=3mm]},thick,warnred] (4,0.6) -- (4,-0.8) node[right] {$F_G$};
\draw[-{Stealth[length=2.5mm]},thick,physikblue] (4,1.4) -- (4,2.2) node[right] {\small $F_W$};
\node at (4,-1.5) {\small $F_G > F_W$};
\node at (4,-2) {\small $\Rightarrow$ beschleunigt noch};

% Phase 3
\node at (8,2.5) {\textbf{Grenzfall}};
\fill[gray!50] (8,1) circle (0.4);
\draw[-{Stealth[length=3mm]},thick,warnred] (8,0.6) -- (8,-0.8) node[right] {$F_G$};
\draw[-{Stealth[length=3mm]},thick,physikblue] (8,1.4) -- (8,2.8) node[right] {\small $F_W$};
\node at (8,-1.5) {\small $F_G = F_W$};
\node at (8,-2) {\small $\Rightarrow$ $v$ konstant!};
\end{tikzpicture}
\end{center}

\begin{formelbox}[Grenzgeschwindigkeit (nur qualitativ!)]
Wenn $F_G = F_W$ gilt, erreicht der Körper seine \textbf{Grenzgeschwindigkeit} -- er fällt dann gleichförmig weiter.

\textbf{Keine Berechnung erforderlich! Nur verstehen:}
\begin{center}
\begin{tabular}{lcl}
Masse $m$ größer & $\Rightarrow$ & $v_{\text{grenz}}$ größer (schwerer = schneller) \\
Fläche $A$ größer & $\Rightarrow$ & $v_{\text{grenz}}$ kleiner (Fallschirm!) \\
Luftdichte $\rho$ größer & $\Rightarrow$ & $v_{\text{grenz}}$ kleiner (mehr Widerstand)
\end{tabular}
\end{center}
\end{formelbox}

\begin{merksatz}
$v_{\text{grenz}}$ wird größer, wenn...
\begin{itemize}
\item die Masse $m$ zunimmt (schwerer $\rightarrow$ schneller)
\item die Fläche $A$ abnimmt (kleiner $\rightarrow$ schneller)
\item die Dichte $\rho$ abnimmt (dünnere Luft $\rightarrow$ schneller)
\item der $c_W$-Wert abnimmt (stromlinienförmiger $\rightarrow$ schneller)
\end{itemize}
\end{merksatz}

\begin{beispiel}
\textbf{Alltagsbeispiele zur Grenzgeschwindigkeit:}

\begin{tabular}{|l|c|l|}
\hline
\rowcolor{lightblue}
\textbf{Objekt} & $v_{\text{grenz}}$ & \textbf{Erklärung} \\
\hline
Regentropfen & $\sim 8$ m/s & Klein, aber auch leicht \\
\hline
Fallschirmspringer (ohne Schirm) & $\sim 50$ m/s & Große Masse, kleine Fläche \\
\hline
Fallschirmspringer (mit Schirm) & $\sim 5$ m/s & \textbf{Große Fläche} bremst! \\
\hline
Feder & $< 1$ m/s & Sehr leicht, große Fläche \\
\hline
\end{tabular}

\vspace{0.3cm}
\textbf{Felix Baumgartner (Red Bull Stratos):}\\
Erreichte $>$1000 km/h -- \textbf{warum so schnell?}\\
$\Rightarrow$ In 39 km Höhe ist die Luft \textbf{300-mal dünner} als am Boden!\\
$\Rightarrow$ Weniger Luftwiderstand $\Rightarrow$ höhere Grenzgeschwindigkeit.
\end{beispiel}

% ============================================================================
% KAPITEL 6: ZUSAMMENFASSUNG
% ============================================================================
\newpage
\section{Formelsammlung}

\subsection{Alle wichtigen Formeln auf einen Blick}

\begin{center}
\renewcommand{\arraystretch}{1.6}
\begin{tabular}{|p{4cm}|c|c|}
\hline
\rowcolor{lightblue}
\textbf{Bewegungsart} & \textbf{Geschwindigkeit} & \textbf{Weg} \\
\hline
\textbf{Gleichförmig} & $v = \text{konst.}$ & $s = v \cdot t + s_0$ \\
\hline
\textbf{Gleichm. beschleunigt} & $v = a \cdot t + v_0$ & $s = \frac{1}{2}at^2 + v_0 t + s_0$ \\
\hline
\textbf{Freier Fall} & $v = g \cdot t$ & $s = \frac{1}{2}g \cdot t^2 = 5t^2$ \\
\hline
\end{tabular}
\end{center}

\vspace{0.5cm}

\begin{center}
\renewcommand{\arraystretch}{1.6}
\begin{tabular}{|l|c|}
\hline
\rowcolor{lightblue}
\textbf{Spezielle Formeln} & \textbf{Formel} \\
\hline
Bremszeit & $t = \dfrac{v_0}{|a|}$ \\
\hline
Bremsweg & $s = \dfrac{v_0^2}{2|a|}$ \\
\hline
Fallzeit (aus Höhe) & $t = \sqrt{\dfrac{2s}{g}} = \sqrt{\dfrac{s}{5}}$ \\
\hline
\hline
\multicolumn{2}{|l|}{\textit{Luftwiderstand: nur qualitativ -- keine Berechnung!}} \\
\hline
\end{tabular}
\end{center}

\subsection{Konstanten und Umrechnungen}

\begin{center}
\begin{tabular}{|l|c|}
\hline
\rowcolor{lightblue}
\textbf{Konstante/Umrechnung} & \textbf{Wert} \\
\hline
Fallbeschleunigung & $g = 9{,}81\,\mathrm{m/s^2} \approx 10\,\mathrm{m/s^2}$ \\
\hline
Luftdichte (Boden) & $\rho = 1{,}2\,\mathrm{kg/m^3}$ \\
\hline
km/h $\rightarrow$ m/s & $\div 3{,}6$ \\
\hline
m/s $\rightarrow$ km/h & $\times 3{,}6$ \\
\hline
\end{tabular}
\end{center}

\subsection{Kopfrechenbare Werte}

\begin{tipp}
\textbf{Diese Werte sicher beherrschen:}
\begin{multicols}{2}
\textbf{Geschwindigkeiten:}
\begin{itemize}
\item 36 km/h = 10 m/s
\item 72 km/h = 20 m/s
\item 90 km/h = 25 m/s
\item 108 km/h = 30 m/s
\end{itemize}

\columnbreak

\textbf{Wurzeln:}
\begin{itemize}
\item $\sqrt{4} = 2$, $\sqrt{9} = 3$
\item $\sqrt{16} = 4$, $\sqrt{25} = 5$
\item $\sqrt{36} = 6$, $\sqrt{49} = 7$
\item $\sqrt{64} = 8$, $\sqrt{81} = 9$, $\sqrt{100} = 10$
\end{itemize}
\end{multicols}
\end{tipp}

% ============================================================================
% KAPITEL 7: ÜBUNGSAUFGABEN MIT LÖSUNGEN
% ============================================================================
\newpage
\section{Übungsaufgaben mit Lösungen}

\begin{aufgabe}[Aufgabe 1: Gleichförmige Bewegung]
Ein ICE fährt mit konstanter Geschwindigkeit 288 km/h.

\textbf{a)} Rechne die Geschwindigkeit in m/s um.

\textbf{b)} Welche Strecke legt der ICE in 30 Sekunden zurück?

\textbf{c)} Wie lange braucht er für 2 km?
\end{aufgabe}

\textbf{Lösung:}
\begin{itemize}
\item[a)] $v = \frac{288}{3{,}6} = 80\,\mathrm{m/s}$
\item[b)] $s = v \cdot t = 80 \cdot 30 = 2400\,\mathrm{m} = 2{,}4\,\mathrm{km}$
\item[c)] $t = \frac{s}{v} = \frac{2000\,\mathrm{m}}{80\,\mathrm{m/s}} = 25\,\mathrm{s}$
\end{itemize}

\vspace{0.5cm}

\begin{aufgabe}[Aufgabe 2: Beschleunigte Bewegung]
Ein Motorrad beschleunigt aus dem Stand mit $a = 5\,\mathrm{m/s^2}$.

\textbf{a)} Welche Geschwindigkeit erreicht es nach 4 Sekunden?

\textbf{b)} Welchen Weg legt es dabei zurück?

\textbf{c)} Welche Geschwindigkeit hat es nach 100 m?
\end{aufgabe}

\textbf{Lösung:}
\begin{itemize}
\item[a)] $v = a \cdot t = 5 \cdot 4 = 20\,\mathrm{m/s} = 72\,\mathrm{km/h}$
\item[b)] $s = \frac{1}{2}at^2 = \frac{1}{2} \cdot 5 \cdot 16 = 40\,\mathrm{m}$
\item[c)] $v^2 = 2as = 2 \cdot 5 \cdot 100 = 1000 \Rightarrow v = \sqrt{1000} \approx 31{,}6\,\mathrm{m/s}$
\end{itemize}

\vspace{0.5cm}

\begin{aufgabe}[Aufgabe 3: Bremsvorgang]
Ein Auto fährt mit 72 km/h und bremst mit $a = -4\,\mathrm{m/s^2}$.

\textbf{a)} Berechne die Bremszeit.

\textbf{b)} Berechne den Bremsweg.
\end{aufgabe}

\textbf{Lösung:}
\begin{itemize}
\item[a)] $v_0 = 72\,\mathrm{km/h} = 20\,\mathrm{m/s}$\\
$t = \frac{v_0}{|a|} = \frac{20}{4} = 5\,\mathrm{s}$
\item[b)] $s = \frac{v_0^2}{2|a|} = \frac{400}{8} = 50\,\mathrm{m}$
\end{itemize}

\vspace{0.5cm}

\begin{aufgabe}[Aufgabe 4: Freier Fall]
Ein Stein wird aus 80 m Höhe fallen gelassen.

\textbf{a)} Berechne die Fallzeit.

\textbf{b)} Bestimme die Aufprallgeschwindigkeit.
\end{aufgabe}

\textbf{Lösung:}
\begin{itemize}
\item[a)] $t = \sqrt{\frac{s}{5}} = \sqrt{\frac{80}{5}} = \sqrt{16} = 4\,\mathrm{s}$
\item[b)] $v = 10 \cdot t = 10 \cdot 4 = 40\,\mathrm{m/s} = 144\,\mathrm{km/h}$
\end{itemize}

\vspace{0.5cm}

\begin{aufgabe}[Aufgabe 5: Luftwiderstand verstehen (qualitativ)]
\textbf{a)} Ein Fallschirmspringer öffnet seinen Fallschirm. Erkläre, warum er danach langsamer fällt.

\textbf{b)} Warum erreichte Felix Baumgartner in 39 km Höhe Überschallgeschwindigkeit, obwohl er in Meereshöhe „nur" etwa 50 m/s schnell wäre?

\textbf{c)} Zwei Kugeln aus dem gleichen Material -- eine mit doppeltem Durchmesser. Welche hat die höhere Grenzgeschwindigkeit? Begründe.
\end{aufgabe}

\textbf{Lösungen:}
\begin{itemize}
\item[a)] Der Fallschirm vergrößert die Querschnittsfläche $A$ enorm $\Rightarrow$ größerer Luftwiderstand $\Rightarrow$ kleinere Grenzgeschwindigkeit.
\item[b)] In großer Höhe ist die Luftdichte $\rho$ viel kleiner $\Rightarrow$ weniger Luftwiderstand $\Rightarrow$ höhere Grenzgeschwindigkeit.
\item[c)] Die größere Kugel. Sie hat 8-mal so viel Masse (Volumen $\sim r^3$), aber nur 4-mal so viel Querschnitt ($A \sim r^2$). Mehr Masse pro Fläche = höhere Grenzgeschwindigkeit.
\end{itemize}

\vspace{0.5cm}

\begin{aufgabe}[Aufgabe 6: Zusammengesetzte Bewegung]
Ein Radfahrer fährt 4 Sekunden lang gleichförmig mit 5 m/s (Startort: $s_0 = 10$ m). Dann bremst er mit $a = -1\,\mathrm{m/s^2}$ bis zum Stillstand.

\textbf{a)} Wo befindet er sich nach 4 s?

\textbf{b)} Wie lange dauert der Bremsvorgang?

\textbf{c)} Welchen Gesamtweg legt er zurück?
\end{aufgabe}

\textbf{Lösung:}
\begin{itemize}
\item[a)] $s(4) = v_0 \cdot t + s_0 = 5 \cdot 4 + 10 = 30\,\mathrm{m}$
\item[b)] $t_{\text{brems}} = \frac{v_0}{|a|} = \frac{5}{1} = 5\,\mathrm{s}$
\item[c)] $s_{\text{brems}} = \frac{v_0^2}{2|a|} = \frac{25}{2} = 12{,}5\,\mathrm{m}$\\
$s_{\text{gesamt}} = 30 + 12{,}5 = 42{,}5\,\mathrm{m}$
\end{itemize}

\vspace{0.8cm}

% ============================================================================
% FINDE DEN FEHLER - AUFGABEN
% ============================================================================
\begin{tcolorbox}[colback=warnred!5, colframe=warnred, boxrule=1pt, title={\textbf{{$\bigstar\bigstar\bigstar$} Finde den Fehler! (AFB III)}}]
Die folgenden Schülerrechnungen enthalten \textbf{typische Fehler}. Finde und korrigiere sie!

\vspace{0.3cm}
\textbf{Fehleraufgabe 1:}\\
\textit{„Ein Auto fährt 72 km/h und bremst 5 s lang mit $a = -2\,\mathrm{m/s^2}$. Wie weit kommt es?"}

\textbf{Schülerrechnung:}
\[
s = \frac{1}{2} \cdot 2 \cdot 5^2 = 25\,\mathrm{m}
\]

\textbf{Was ist falsch?}\\[2.5cm]

\textbf{Fehleraufgabe 2:}\\
\textit{„Ein Stein fällt 3 Sekunden. Mit welcher Geschwindigkeit kommt er unten an?"}

\textbf{Schülerrechnung:}
\[
v = g \cdot t = 9{,}81 \cdot 3 = 29{,}43\,\mathrm{m/s^2}
\]

\textbf{Was ist falsch?}\\[2cm]
\end{tcolorbox}

\textbf{Lösungen Fehleraufgaben:}\\
\textbf{1.} Zwei Fehler: (a) Der Schüler hat die Anfangsgeschwindigkeit $v_0 = 20\,\mathrm{m/s}$ vergessen! (b) Der Betrag von $a$ wurde verwendet, ohne das Vorzeichen zu berücksichtigen. Richtig: $s = v_0 \cdot t + \frac{1}{2}at^2 = 20 \cdot 5 + \frac{1}{2}(-2) \cdot 25 = 100 - 25 = 75\,\mathrm{m}$\\
\textbf{2.} Zwei Fehler: (a) Die Einheit m/s² ist falsch -- Geschwindigkeit hat die Einheit \textbf{m/s}. (b) Wir verwenden $g \approx 10\,\mathrm{m/s^2}$ für Kopfrechnen! Richtig: $v = 10 \cdot 3 = 30\,\mathrm{m/s}$

% ============================================================================
% KAPITEL 8: SELBSTTEST
% ============================================================================
\newpage
\section{Selbsttest: Habe ich alles verstanden?}

\subsection{Lückentext -- Fülle die Lücken aus!}

\begin{tcolorbox}[colback=lightgray, colframe=gray, boxrule=0.5pt]
\begin{enumerate}[itemsep=8pt]
\item Bei einer \underline{\hspace{3cm}} Bewegung bleibt die Geschwindigkeit konstant.

\item Im $s(t)$-Diagramm entspricht die \underline{\hspace{3cm}} der Geschwindigkeit.

\item Die Beschleunigung gibt an, wie stark sich die \underline{\hspace{3cm}} pro Zeiteinheit ändert.

\item Beim freien Fall gilt $g \approx$ \underline{\hspace{2cm}} m/s².

\item Die Formel für die Fallstrecke (vereinfacht) lautet: $s = $ \underline{\hspace{2cm}}.

\item Die Luftwiderstandskraft hängt von $v^{\underline{\hspace{1cm}}}$ ab.

\item Wenn $F_G = F_W$ gilt, erreicht ein fallender Körper seine \underline{\hspace{3cm}}.

\item 72 km/h entspricht \underline{\hspace{2cm}} m/s.
\end{enumerate}
\end{tcolorbox}

\textbf{Lösungen:} 1) gleichförmigen, 2) Steigung, 3) Geschwindigkeit, 4) 10, 5) $5t^2$, 6) 2, 7) Grenzgeschwindigkeit, 8) 20

\subsection{Quick-Check: Richtig oder Falsch?}

\begin{center}
\begin{tabular}{|p{10cm}|c|c|}
\hline
\rowcolor{lightblue}
\textbf{Aussage} & \textbf{R} & \textbf{F} \\
\hline
Im $v(t)$-Diagramm einer gleichförmigen Bewegung ist der Graph eine horizontale Linie. & $\square$ & $\square$ \\
\hline
Beim Bremsen ist die Beschleunigung positiv. & $\square$ & $\square$ \\
\hline
Im freien Fall erreichen alle Körper (ohne Luftwiderstand) die gleiche Geschwindigkeit. & $\square$ & $\square$ \\
\hline
Je größer die Querschnittsfläche, desto größer die Grenzgeschwindigkeit. & $\square$ & $\square$ \\
\hline
Die Fläche unter dem $v(t)$-Graphen entspricht dem zurückgelegten Weg. & $\square$ & $\square$ \\
\hline
\end{tabular}
\end{center}

\textbf{Lösungen:} R, F, R, F, R

% ============================================================================
% KAPITEL 9: DIFFERENZIERTE ÜBUNGEN
% ============================================================================
\newpage
\section{Differenzierte Übungsaufgaben}

\begin{tcolorbox}[colback=merkegreen!10, colframe=merkegreen, boxrule=1pt]
\textbf{Schwierigkeitsgrade:}
\begin{itemize}
\item[{$\bigstar$}] Grundlagen -- Einsetzen und berechnen
\item[{$\bigstar\bigstar$}] Fortgeschritten -- Umstellen und mehrere Schritte
\item[{$\bigstar\bigstar\bigstar$}] Anspruchsvoll -- Transfer und Begründen
\end{itemize}
\end{tcolorbox}

\vspace{0.5cm}

\begin{aufgabe}[{$\bigstar$} Aufgabe A: Grundlagen]
\textbf{a)} Rechne um: 54 km/h = \underline{\hspace{2cm}} m/s

\textbf{b)} Ein Stein fällt 2 Sekunden. Berechne die Fallstrecke.

\textbf{c)} Ein Auto fährt 90 km/h. Welchen Weg legt es in 10 s zurück?
\end{aufgabe}

\textbf{Lösungen:} a) 15 m/s \quad b) $s = 5 \cdot 4 = 20$ m \quad c) $s = 25 \cdot 10 = 250$ m

\vspace{0.5cm}

\begin{aufgabe}[{$\bigstar\bigstar$} Aufgabe B: Fortgeschritten]
\textbf{a)} Ein Auto beschleunigt von 0 auf 72 km/h in 8 Sekunden. Berechne die Beschleunigung.

\textbf{b)} Ein Stein trifft mit 30 m/s auf dem Boden auf. Aus welcher Höhe wurde er fallen gelassen?

\textbf{c)} Zwei Autos starten gleichzeitig. Auto A: $s_0 = 100$ m, $v_A = 15$ m/s. Auto B: $s_0 = 0$, $v_B = 20$ m/s. Wann und wo überholt B das Auto A?
\end{aufgabe}

\textbf{Lösungen:}\\
a) $v = 72\,\mathrm{km/h} = 20\,\mathrm{m/s}$, $a = \frac{20}{8} = 2{,}5\,\mathrm{m/s^2}$\\
b) $v = 10t \Rightarrow t = 3\,\mathrm{s}$, $s = 5 \cdot 9 = 45\,\mathrm{m}$\\
c) $15t + 100 = 20t \Rightarrow t = 20\,\mathrm{s}$, $s = 400\,\mathrm{m}$

\vspace{0.5cm}

\begin{aufgabe}[{$\bigstar\bigstar\bigstar$} Aufgabe C: Transfer und Begründen (AFB III)]
\textbf{a)} Ein Fallschirmspringer springt aus 4000 m Höhe. Erkläre, warum er \textbf{nicht} mit der berechneten Geschwindigkeit aus $v = \sqrt{2gs}$ auf dem Boden aufkommt.

\textbf{b)} Beurteile: Ein Schüler behauptet: „Im Vakuum fallen schwere Gegenstände schneller als leichte." Ist das korrekt? Begründe physikalisch.

\textbf{c)} Ein Konstrukteur möchte die Grenzgeschwindigkeit eines Fallschirms halbieren, ohne die Masse zu ändern. Welche Möglichkeiten hat er? Begründe qualitativ.
\end{aufgabe}

\textbf{Lösungen:}\\
a) Der Luftwiderstand begrenzt die Geschwindigkeit auf die Grenzgeschwindigkeit ($\approx 50$ m/s ohne Schirm, $\approx 5$ m/s mit Schirm). Die Formel $v = \sqrt{2gs}$ gilt nur im Vakuum.\\
b) Falsch. Im Vakuum gilt $s = \frac{1}{2}gt^2$ -- die Masse kommt nicht vor. Alle Körper fallen gleich schnell (Galilei).\\
c) Er muss den Luftwiderstand erhöhen: größere Schirmfläche $A$ oder höherer Widerstandsbeiwert $c_W$ (andere Schirmform). Mehr Luftwiderstand = niedrigere Grenzgeschwindigkeit.

\vspace{0.5cm}

\begin{aufgabe}[{$\bigstar\bigstar\bigstar$} Aufgabe D: Komplexe Analyse]
Das folgende $v(t)$-Diagramm zeigt die Bewegung eines Radfahrers:

\begin{center}
\begin{tikzpicture}
\begin{axis}[
    width=10cm, height=5cm,
    xlabel={$t$ in s},
    ylabel={$v$ in m/s},
    xmin=0, xmax=12,
    ymin=0, ymax=12,
    grid=major,
    axis lines=left,
    xtick={0,2,4,6,8,10,12},
    ytick={0,2,4,6,8,10,12}
]
\addplot[thick, physikblue, mark=*] coordinates {(0,0) (4,8) (8,8) (10,0)};
\end{axis}
\end{tikzpicture}
\end{center}

\textbf{a)} Beschreibe die Bewegung in Worten.

\textbf{b)} Berechne die Beschleunigung in den ersten 4 Sekunden.

\textbf{c)} Berechne die Bremsverzögerung am Ende.

\textbf{d)} Berechne den Gesamtweg (Tipp: Fläche unter dem Graphen).

\textbf{e)} Skizziere das zugehörige $s(t)$-Diagramm.
\end{aufgabe}

\textbf{Lösungen:}\\
a) Beschleunigung (0-4s), gleichförmig (4-8s), Bremsen (8-10s)\\
b) $a = \frac{8}{4} = 2\,\mathrm{m/s^2}$\\
c) $a = \frac{0-8}{2} = -4\,\mathrm{m/s^2}$\\
d) Dreieck + Rechteck + Dreieck = $\frac{1}{2} \cdot 4 \cdot 8 + 4 \cdot 8 + \frac{1}{2} \cdot 2 \cdot 8 = 16 + 32 + 8 = 56\,\mathrm{m}$\\
e) Parabel (0-4s), Gerade (4-8s), nach unten gekrümmte Kurve (8-10s)

% ============================================================================
% KOMPLEXE DIAGRAMM-ÜBUNGEN
% ============================================================================
\newpage
\subsection{Komplexe Diagramm-Aufgaben zum Interpretieren}

\begin{tcolorbox}[colback=physikblue!5, colframe=physikblue, boxrule=1pt]
\textbf{Hinweis:} Diese Aufgaben trainieren das \textbf{Lesen und Interpretieren} von Bewegungsdiagrammen -- eine typische Klausuraufgabe! Nimm dir Zeit und arbeite systematisch.
\end{tcolorbox}

\vspace{0.5cm}

\begin{aufgabe}[{$\bigstar\bigstar\bigstar$} Diagramm E: Komplexes s-t-Diagramm]
Das folgende \textbf{Weg-Zeit-Diagramm} zeigt die Bewegung eines Autos auf einer Landstraße:

\begin{center}
\begin{tikzpicture}
\begin{axis}[
    width=12cm, height=6cm,
    xlabel={$t$ in s},
    ylabel={$s$ in m},
    xmin=0, xmax=25,
    ymin=0, ymax=350,
    grid=major,
    axis lines=left,
    xtick={0,5,10,15,20,25},
    ytick={0,50,100,150,200,250,300,350},
    legend pos=north west
]
% Phase 1: Beschleunigung (0-5s), Parabel s = 2.5t²
\addplot[thick, physikblue, domain=0:5, samples=50] {2.5*x^2};
% Phase 2: Gleichförmig (5-15s), v=25 m/s, von s(5)=62.5 weiter
\addplot[thick, physikblue, domain=5:15] {62.5 + 25*(x-5)};
% Phase 3: Bremsen (15-20s), von s=312.5, v=25 auf 0
\addplot[thick, physikblue, domain=15:20, samples=50] {312.5 + 25*(x-15) - 2.5*(x-15)^2};
% Phase 4: Stillstand (20-25s)
\addplot[thick, physikblue, domain=20:25] {375};
% Punkte markieren
\addplot[only marks, mark=*, mark size=2pt, physikblue] coordinates {(0,0) (5,62.5) (15,312.5) (20,375) (25,375)};
\node at (axis cs:2.5,40) {\small I};
\node at (axis cs:10,200) {\small II};
\node at (axis cs:17.5,350) {\small III};
\node at (axis cs:22.5,355) {\small IV};
\end{axis}
\end{tikzpicture}
\end{center}

\textbf{Aufgaben:}
\begin{enumerate}[label=\alph*)]
\item Beschreibe die vier Phasen I--IV in Worten.
\item Lies die Geschwindigkeit in Phase II aus dem Diagramm ab.
\item Berechne die Beschleunigung in Phase I.
\item Berechne die Verzögerung in Phase III.
\item Welchen Gesamtweg legt das Auto zurück?
\item Zeichne das zugehörige $v(t)$-Diagramm.
\end{enumerate}
\end{aufgabe}

\vspace{0.5cm}

\begin{aufgabe}[{$\bigstar\bigstar\bigstar$} Diagramm F: Zwei Fahrzeuge im v-t-Diagramm]
Das \textbf{Geschwindigkeit-Zeit-Diagramm} zeigt zwei Fahrzeuge, die zur gleichen Zeit am gleichen Ort starten:

\begin{center}
\begin{tikzpicture}
\begin{axis}[
    width=12cm, height=6cm,
    xlabel={$t$ in s},
    ylabel={$v$ in m/s},
    xmin=0, xmax=12,
    ymin=0, ymax=25,
    grid=major,
    axis lines=left,
    xtick={0,2,4,6,8,10,12},
    ytick={0,5,10,15,20,25},
    legend pos=north east
]
% Fahrzeug A: Beschleunigt auf 20 m/s in 4s, dann konstant
\addplot[thick, physikblue, mark=none] coordinates {(0,0) (4,20) (12,20)};
% Fahrzeug B: Konstant 15 m/s von Anfang an
\addplot[thick, merkegreen, dashed, mark=none] coordinates {(0,15) (12,15)};
\legend{Fahrzeug A, Fahrzeug B}
\end{axis}
\end{tikzpicture}
\end{center}

\textbf{Aufgaben:}
\begin{enumerate}[label=\alph*)]
\item Beschreibe die Bewegung von Fahrzeug A und Fahrzeug B.
\item Berechne die Beschleunigung von Fahrzeug A in den ersten 4 Sekunden.
\item Berechne den Weg, den Fahrzeug A in den ersten 4 Sekunden zurücklegt.
\item Berechne den Weg, den Fahrzeug B in den ersten 4 Sekunden zurücklegt.
\item Welches Fahrzeug ist nach 4 s vorne? Um wie viel?
\item Nach welcher Zeit haben beide Fahrzeuge den gleichen Weg zurückgelegt?\\
\textit{Tipp: Stelle die Weg-Gleichungen auf und setze sie gleich.}
\end{enumerate}
\end{aufgabe}

\vspace{0.5cm}

\begin{aufgabe}[{$\bigstar\bigstar\bigstar$} Diagramm G: Vom v-t- zum s-t-Diagramm]
Gegeben ist folgendes \textbf{v-t-Diagramm}:

\begin{center}
\begin{tikzpicture}
\begin{axis}[
    width=12cm, height=5cm,
    xlabel={$t$ in s},
    ylabel={$v$ in m/s},
    xmin=0, xmax=16,
    ymin=-10, ymax=15,
    grid=major,
    axis lines=center,
    xtick={0,2,4,6,8,10,12,14,16},
    ytick={-10,-5,0,5,10,15}
]
% Phase 1: Konstant 10 m/s (0-4s)
\addplot[thick, physikblue] coordinates {(0,10) (4,10)};
% Phase 2: Abbremsen auf 0 (4-8s)
\addplot[thick, physikblue] coordinates {(4,10) (8,0)};
% Phase 3: Rückwärts beschleunigen auf -5 m/s (8-10s)
\addplot[thick, physikblue] coordinates {(8,0) (10,-5)};
% Phase 4: Konstant -5 m/s (10-14s)
\addplot[thick, physikblue] coordinates {(10,-5) (14,-5)};
% Phase 5: Abbremsen auf 0 (14-16s)
\addplot[thick, physikblue] coordinates {(14,-5) (16,0)};
% Markierungen
\addplot[only marks, mark=*, mark size=2pt, physikblue] coordinates {(0,10) (4,10) (8,0) (10,-5) (14,-5) (16,0)};
\end{axis}
\end{tikzpicture}
\end{center}

\textbf{Aufgaben:}
\begin{enumerate}[label=\alph*)]
\item Was bedeutet eine \textbf{negative Geschwindigkeit}?
\item Berechne die Beschleunigung in jeder Phase.
\item Berechne den zurückgelegten Weg für jede Phase (Fläche!). Beachte die Vorzeichen!
\item An welchem Ort (bezogen auf den Start) befindet sich das Objekt nach 16 s?
\item Skizziere das zugehörige $s(t)$-Diagramm.
\end{enumerate}
\end{aufgabe}

\newpage
\textbf{Lösungen zu den komplexen Diagramm-Aufgaben:}

\vspace{0.3cm}
\textbf{Diagramm E (s-t-Diagramm):}
\begin{itemize}
\item[a)] \textbf{Phase I:} Beschleunigte Bewegung (Parabel nach oben). \textbf{Phase II:} Gleichförmige Bewegung (Gerade mit konstantem Anstieg). \textbf{Phase III:} Abbremsen (Parabel flacht ab). \textbf{Phase IV:} Stillstand (horizontale Linie).
\item[b)] Steigung in Phase II: $v = \frac{\Delta s}{\Delta t} = \frac{312{,}5 - 62{,}5}{15-5} = \frac{250}{10} = \boxed{25\,\mathrm{m/s}}$
\item[c)] In Phase I gilt $s = \frac{1}{2}at^2$. Mit $s(5) = 62{,}5$ m und $t = 5$ s:\\
$62{,}5 = \frac{1}{2} \cdot a \cdot 25 \Rightarrow a = \frac{62{,}5 \cdot 2}{25} = \boxed{5\,\mathrm{m/s^2}}$
\item[d)] In Phase III bremst das Auto von 25 m/s auf 0 in 5 s:\\
$a = \frac{0 - 25}{5} = \boxed{-5\,\mathrm{m/s^2}}$
\item[e)] Gesamtweg: $s_{\text{ges}} = \boxed{375\,\mathrm{m}}$ (ablesbar als Endposition)
\item[f)] v-t-Diagramm: Gerade (0→25 m/s in 5s), Horizontale (25 m/s für 10s), Gerade (25→0 m/s in 5s), Horizontale bei 0.
\end{itemize}

\vspace{0.3cm}
\textbf{Diagramm F (Zwei Fahrzeuge):}
\begin{itemize}
\item[a)] \textbf{A:} Beschleunigt 4s, dann gleichförmig mit 20 m/s. \textbf{B:} Fährt die ganze Zeit gleichförmig mit 15 m/s.
\item[b)] $a_A = \frac{20}{4} = \boxed{5\,\mathrm{m/s^2}}$
\item[c)] Fläche unter dem Dreieck: $s_A(4) = \frac{1}{2} \cdot 4 \cdot 20 = \boxed{40\,\mathrm{m}}$
\item[d)] Rechteck: $s_B(4) = 15 \cdot 4 = \boxed{60\,\mathrm{m}}$
\item[e)] \textbf{Fahrzeug B ist vorne}, um $60 - 40 = \boxed{20\,\mathrm{m}}$
\item[f)] Weggleichungen: $s_A(t) = \frac{1}{2} \cdot 5 \cdot t^2 = 2{,}5t^2$ (für $t \leq 4$s) und ab $t > 4$: $s_A(t) = 40 + 20(t-4)$.\\
$s_B(t) = 15t$.\\
Gleichsetzen für $t > 4$: $40 + 20(t-4) = 15t \Rightarrow 40 + 20t - 80 = 15t \Rightarrow 5t = 40 \Rightarrow \boxed{t = 8\,\mathrm{s}}$\\
Kontrolle: $s_A(8) = 40 + 20 \cdot 4 = 120$ m, $s_B(8) = 15 \cdot 8 = 120$ m \checkmark
\end{itemize}

\vspace{0.3cm}
\textbf{Diagramm G (Negative Geschwindigkeit):}
\begin{itemize}
\item[a)] Negative Geschwindigkeit bedeutet \textbf{Bewegung in Gegenrichtung} (z.\,B. rückwärts fahren).
\item[b)] Beschleunigungen:\\
Phase 1 (0-4s): $a = 0$ (konstant)\\
Phase 2 (4-8s): $a = \frac{0-10}{4} = -2{,}5\,\mathrm{m/s^2}$ (Bremsen)\\
Phase 3 (8-10s): $a = \frac{-5-0}{2} = -2{,}5\,\mathrm{m/s^2}$ (Rückwärts beschleunigen)\\
Phase 4 (10-14s): $a = 0$ (konstant rückwärts)\\
Phase 5 (14-16s): $a = \frac{0-(-5)}{2} = +2{,}5\,\mathrm{m/s^2}$ (Abbremsen der Rückwärtsbewegung)
\item[c)] Wege (Flächen):\\
Phase 1: $s_1 = 10 \cdot 4 = +40$ m (Rechteck)\\
Phase 2: $s_2 = \frac{1}{2} \cdot 4 \cdot 10 = +20$ m (Dreieck)\\
Phase 3: $s_3 = \frac{1}{2} \cdot 2 \cdot (-5) = -5$ m (Dreieck unter x-Achse)\\
Phase 4: $s_4 = (-5) \cdot 4 = -20$ m (Rechteck unter x-Achse)\\
Phase 5: $s_5 = \frac{1}{2} \cdot 2 \cdot (-5) = -5$ m (Dreieck unter x-Achse)
\item[d)] Endposition: $s_{\text{ges}} = 40 + 20 - 5 - 20 - 5 = \boxed{+30\,\mathrm{m}}$\\
Das Objekt ist 30 m \textbf{vor} dem Startpunkt (in Vorwärtsrichtung).
\item[e)] s-t-Diagramm: Gerade (0-4s, Steigung +10), nach oben gekrümmte Parabel bis zum Maximum bei t=8s, dann Abstieg (erst Parabel, dann Gerade, dann abflachend) bis s=30m bei t=16s.
\end{itemize}

% ============================================================================
% CHECKLISTE
% ============================================================================
\newpage
\section{Checkliste zur Klausurvorbereitung}

\begin{tcolorbox}[colback=lightblue, colframe=physikblue, title=\textbf{Kann ich das? -- Ehrlich ankreuzen!}]
\begin{tabular}{|p{9cm}|c|c|c|}
\hline
\rowcolor{lightblue}
\textbf{Kompetenz} & \textbf{Ja} & \textbf{Teils} & \textbf{Nein} \\
\hline
\multicolumn{4}{|l|}{\textbf{Gleichförmige Bewegung}} \\
\hline
km/h in m/s umrechnen und umgekehrt & $\square$ & $\square$ & $\square$ \\
\hline
$s(t)$- und $v(t)$-Diagramme zeichnen & $\square$ & $\square$ & $\square$ \\
\hline
Steigung im $s(t)$-Diagramm interpretieren & $\square$ & $\square$ & $\square$ \\
\hline
\multicolumn{4}{|l|}{\textbf{Beschleunigte Bewegung}} \\
\hline
Formeln $v(t)$ und $s(t)$ anwenden & $\square$ & $\square$ & $\square$ \\
\hline
Bremsweg und Bremszeit berechnen & $\square$ & $\square$ & $\square$ \\
\hline
Anfangsbedingungen $v_0$ und $s_0$ einsetzen & $\square$ & $\square$ & $\square$ \\
\hline
\multicolumn{4}{|l|}{\textbf{Freier Fall}} \\
\hline
Fallzeit aus Höhe berechnen & $\square$ & $\square$ & $\square$ \\
\hline
Aufprallgeschwindigkeit bestimmen & $\square$ & $\square$ & $\square$ \\
\hline
Senkrechten Wurf nach oben lösen & $\square$ & $\square$ & $\square$ \\
\hline
\multicolumn{4}{|l|}{\textbf{Luftwiderstand \& Grenzgeschwindigkeit} \textit{(nur qualitativ!)}} \\
\hline
Erklären, warum $F_W$ bei höherer Geschwindigkeit zunimmt & $\square$ & $\square$ & $\square$ \\
\hline
Erklären, wann Grenzgeschwindigkeit erreicht wird ($F_G = F_W$) & $\square$ & $\square$ & $\square$ \\
\hline
Einfluss von Masse, Fläche, Luftdichte auf $v_{\text{grenz}}$ beschreiben & $\square$ & $\square$ & $\square$ \\
\hline
\end{tabular}
\end{tcolorbox}

\vspace{0.5cm}

\begin{warnung}
\textbf{Die häufigsten Fehler in Klausuren:}
\begin{enumerate}
\item \textbf{Einheiten vergessen} -- Immer km/h in m/s umrechnen!
\item \textbf{$g = 9{,}81$ statt $g \approx 10$} -- Kopfrechnen unmöglich!
\item \textbf{Anfangsbedingungen vergessen} -- $s_0$ und $v_0$ prüfen!
\item \textbf{Vorzeichen bei Bremsung} -- $a < 0$ bei Verzögerung!
\item \textbf{Fläche vs. Steigung verwechselt} -- Im $v(t)$-Diagramm: Fläche = Weg!
\end{enumerate}
\end{warnung}

\vspace{1cm}

\begin{center}
\large\textbf{Viel Erfolg bei der Klausur!}
\end{center}

% ============================================================================
% ANHANG: HERLEITUNG DER ZEITLOSEN FORMEL
% ============================================================================
\newpage
\appendix
\section{Herleitung der zeitlosen Formel}
\label{sec:herleitung}

\begin{tcolorbox}[colback=physikblue!5, colframe=physikblue, boxrule=1pt, title=\textbf{Für Interessierte: Woher kommt $v^2 = v_0^2 + 2as$?}]
Diese Formel ist keine neue Physik, sondern entsteht durch \textbf{geschicktes Kombinieren} der beiden Grundformeln. Der Trick: Wir \textbf{eliminieren die Zeit $t$}.
\end{tcolorbox}

\vspace{0.5cm}

\textbf{Ausgangspunkt: Die zwei Grundformeln}

Wir starten mit den bekannten Formeln für gleichmäßig beschleunigte Bewegung:

\begin{align}
v(t) &= v_0 + a \cdot t \tag{1: Geschwindigkeit}\\[0.5em]
s(t) &= s_0 + v_0 \cdot t + \frac{1}{2} a \cdot t^2 \tag{2: Weg}
\end{align}

Zur Vereinfachung setzen wir $s_0 = 0$ (Start am Ursprung).

\vspace{0.5cm}
\textbf{Schritt 1: Formel (1) nach $t$ umstellen}

\begin{align*}
v &= v_0 + a \cdot t \\[0.3em]
v - v_0 &= a \cdot t \\[0.3em]
t &= \frac{v - v_0}{a}
\end{align*}

\vspace{0.5cm}
\textbf{Schritt 2: Dieses $t$ in Formel (2) einsetzen}

\[
s = v_0 \cdot \underbrace{\frac{v - v_0}{a}}_{= t} + \frac{1}{2} a \cdot \underbrace{\left(\frac{v - v_0}{a}\right)^2}_{= t^2}
\]

\vspace{0.5cm}
\textbf{Schritt 3: Ausmultiplizieren}

Erster Term:
\[
v_0 \cdot \frac{v - v_0}{a} = \frac{v_0 \cdot v - v_0^2}{a}
\]

Zweiter Term (mit $t^2 = \frac{(v-v_0)^2}{a^2}$):
\[
\frac{1}{2} a \cdot \frac{(v - v_0)^2}{a^2} = \frac{(v - v_0)^2}{2a} = \frac{v^2 - 2v \cdot v_0 + v_0^2}{2a}
\]

\vspace{0.5cm}
\textbf{Schritt 4: Zusammenfassen}

\begin{align*}
s &= \frac{v_0 \cdot v - v_0^2}{a} + \frac{v^2 - 2v \cdot v_0 + v_0^2}{2a}
\end{align*}

Auf gemeinsamen Nenner $2a$ bringen:
\begin{align*}
s &= \frac{2(v_0 \cdot v - v_0^2)}{2a} + \frac{v^2 - 2v \cdot v_0 + v_0^2}{2a} \\[0.5em]
s &= \frac{2v_0 v - 2v_0^2 + v^2 - 2v_0 v + v_0^2}{2a} \\[0.5em]
s &= \frac{v^2 - v_0^2}{2a}
\end{align*}

\vspace{0.5cm}
\textbf{Schritt 5: Nach $v^2$ umstellen}

\begin{align*}
s \cdot 2a &= v^2 - v_0^2 \\[0.3em]
v^2 &= v_0^2 + 2 \cdot a \cdot s
\end{align*}

\begin{formelbox}[Ergebnis: Die zeitlose Formel]
\[
\boxed{v^2 = v_0^2 + 2 \cdot a \cdot s}
\]
Diese Formel verbindet \textbf{Geschwindigkeit} und \textbf{Strecke} direkt -- ohne den Umweg über die Zeit!
\end{formelbox}

\vspace{0.5cm}
\begin{merksatz}
\textbf{Wann ist die zeitlose Formel praktisch?}
\begin{itemize}
\item Wenn die \textbf{Zeit nicht gegeben} und nicht gefragt ist
\item Bei \textbf{Bremsweg-Berechnungen} (Endgeschwindigkeit $v = 0$)
\item Bei \textbf{Fallproblemen}: $v^2 = 2gh$ (aus Ruhe, mit $g$ statt $a$)
\end{itemize}
\end{merksatz}

\vspace{0.5cm}
\begin{tcolorbox}[colback=tipporange!10, colframe=tipporange, boxrule=0.5pt]
\textbf{Schnelltest:} Kannst du zeigen, dass für den freien Fall (aus Ruhe, $v_0 = 0$, $a = g$):
\[
v = \sqrt{2gh}
\]
gilt? Probier es mit der zeitlosen Formel!
\end{tcolorbox}

\end{document}
